\documentclass{article}
\usepackage{amsmath, amssymb, amsthm}
\usepackage{hyperref}
\usepackage{geometry}
\geometry{margin=1in}

\title{Erd\H{o}s Problem \#306}
\date{Last updated: 2025-08-31}
\author{Source: erdosproblems.com}

\begin{document}
\maketitle

\noindent\textbf{Status:} OPEN \quad \textbf{No prize}

\noindent\textbf{Tags:} number theory, unit fractions

\medskip
\noindent\textbf{Problem Statement:}

Let $a/b\in \mathbb{Q}_{>0}$ with $b$ squarefree. Are there integers $1<n_1<\cdots<n_k$, each the product of two distinct primes, such that\[\frac{a}{b}=\frac{1}{n_1}+\cdots+\frac{1}{n_k}?\]

\bigskip
\noindent\textbf{Remarks:}

For $n_i$ the product of three distinct primes, this is true when $b=1$, as proved by Butler, Erd\H{o}s and Graham \cite{BEG15} (this paper is perhaps Erd\H{o}s' last paper, appearing 19 years after his death).

There are several examples known of such decompositions for $a/b=1$. The first was given by Barbeau \cite{Ba77}. The shortest known examples have $47$ terms, discovered by Watanabe \cite{Wa20} (see the introduction of \cite{Wa20} for further history).

\bigskip
\noindent\textbf{References:}

[BEG15] Butler, Steve and Erd\H{o}s, Paul and Graham, Ron, Egyptian fractions with each denominator having three distinct
prime divisors. Integers (2015), Paper No. A51, 9.
 
 [Ba77] Barbeau, E. J., Expressing one as a sum of distinct reciprocals. Eureka (Ottawa) (1977), 178-181.
 
 [Wa20] T. Watanabe, New examples of the representation of 1 by the sum of reciprocals of semiprime numbers. arXiv:2009.03275 (2020).

\bigskip
\noindent\small{Source: \url{https://www.erdosproblems.com/306}}
\end{document}
